Se observará que si, en el enunciado, se elimina la hipótesis
2\textsuperscript{o} todavía se puede concluir, en virtud de (3.3.1), la relación

\begin{env}[3.3.6.4]
\label{IV.2.3.3.6.4-fr}
\[
\operatorname{Ass}(\mathcal F\otimes_S\mathcal G)
\supset\bigcup_{x\in\operatorname{Ass}(\mathcal F)}
\left(\bigcup_{y\in\operatorname{Ass}(\mathcal G_{f(x)})}I_{x,y}\right).
\]
\end{env}

\begin{corollary}[3.3.7]
\label{IV.2.3.3.7-fr}
Bajo las hipótesis de \textup{(3.3.6)}, supongamos además que $S$ sea localmente
noetheriano y que
$f(\operatorname{Ass}(\mathcal F))\subset\operatorname{Ass}(\mathcal O_S)$.
Entonces se tiene
\end{corollary}

\begin{env}[3.3.7.1]
\label{IV.2.3.3.7.1-fr}
\[
\operatorname{Ass}(\mathcal F\otimes_S\mathcal G)
=\bigcup_{(x,y)\in C}I_{x,y}
\]
\end{env}

\textit{donde $C$ es el conjunto de pares $(x,y)$ tales que
$x\in\operatorname{Ass}(\mathcal F)$, $y\in\operatorname{Ass}(\mathcal G)$ y
$f(x)=g(y)$.}

Como $\mathcal G$ es plano respecto de $g$, se deduce en efecto de (3.3.1) que la relación
«~$s\in\operatorname{Ass}(\mathcal O_S)$ y
$y\in\operatorname{Ass}(\mathcal G_s)$~» equivale a
$y\in\operatorname{Ass}(\mathcal G)$: la conclusión se sigue de (3.3.6.1).

\begin{remarks}[3.3.8]
\label{IV.2.3.3.8-fr}
Veremos más adelante (4.2.2) que, bajo las hipótesis de (3.3.6),
$T_{x,y}$ es un preesquema \emph{sin ciclo primo asociado inmerso}; de ello se
deducirá que si $\mathcal F$ y los $\mathcal G_s$ carecen de ciclo primo
asociado inmerso, lo mismo ocurre con $\mathcal F\otimes_S\mathcal G$.
\end{remarks}

\begin{corollary}[3.3.9]
\label{IV.2.3.3.9-fr}
Bajo las condiciones de \textup{(3.3.7)}, se tiene
\end{corollary}

\begin{env}[3.3.9.1]
\label{IV.2.3.3.9.1-fr}
\[
q(\operatorname{Ass}(\mathcal F\otimes_S\mathcal G))
\subset\operatorname{Ass}(\mathcal G)
\]
\end{env}

\textit{(donde $q:X\times_S Y\to Y$ es la proyección canónica).}

En efecto, si $(x,y)\in Z$, se tiene
$q(I_{x,y})=\{y\}\subset\operatorname{Ass}(\mathcal G)$.
